%% kapa: exam=1 round=37 year=2026 period=1 form=A subject=real_estate subject_ko=부동산학원론 seq=29 number=109
%% source: Q-Net 감정평가사 기출문제 (cst00309) artlSeq=5258924 file=raw/1008/5258924/2026년 제37회 감정평가사 1차 1교시.pdf
%% conversion: kordoc 4.24.51 -> tools/gen1.py (auto); review=proofread-sonnet; flags=none
\begin{question}{109}
한국주택금융공사의 주택담보노후연금(주택연금)에 관한 설명으로 옳은 것은?
\begin{choices}
  \item 가입자가 오래 생존하는 장수명(長壽命)위험은 공사가 부담하며, 가입 이후 담보주택의 가격이 하락하는 가격위험은 가입자가 부담한다.
  \item 다른 조건이 동일한 경우 담보주택가격이 높을수록, 가입자의 나이가 적을수록 월지급금은 높아진다.
  \item 다른 조건이 동일한 경우 주택연금 가입 후 담보주택의 가격이 상승하면 월지급금도 주택가격에 연동되어 상승한다.
  \item 신청인의 보유주택수가 1주택이라면 주택가격과 거주여부에 관계없이 가입이 가능하며, 2주택 이상인 경우 1주택 이외의 주택 처분을 조건으로 가입이 가능하다.
  \item 월지급금의 기준이 되는 주택가격은 공신력 있는 기관의 인터넷 시세 또는 공시가격을 적용하며, 신청인이 원하면 외부 감정평가업자(감정평가법인등)의 감정평가액으로 할 수 있다.
\end{choices}
\end{question}
\answer{5}
